
% Das ist der KIT-Stil
%\usepackage{../TutTexbib/beamerthemekit}
\usepackage[deutsch,titlepage0]{../TutTexbib/KIT/beamerthemeKITmod}
\TitleImage[width=\titleimagewd]{../figures/titlepage.jpg}
%\usetheme[deutsch,titlepage0]{KIT}

% Include PDFs
\usepackage{pdfpages}

% Libertine font (Original GBI font)
\usepackage{libertine}
%\renewcommand*\familydefault{\sfdefault}  %% Only if the base font of the document is to be sans serif

% Nicer math symbols
\usepackage{eulervm}
%\usepackage{mathpazo}

%% Deutsche Silbentrennung und Beschriftungen
\usepackage[ngerman]{babel}



\usepackage{csquotes}

%% Bibliotheken für viele mathematische Symbole
\usepackage{amsmath, amsfonts, amssymb}

%% Anzeigetiefe für Inhaltsverzeichnis: 1 Stufe
\setcounter{tocdepth}{1}

%% Schönere Schriften
\usepackage[TS1,T1]{fontenc}

%% Tabellen
\usepackage{array}
\usepackage{multicol}

%% Bibliothek für Graphiken
\usepackage{graphicx}

%% der wird sowieso in jeder Datei gesetzt
%%\graphicspath{{../figures/}}

%% Hyperlinks
\usepackage{hyperref}
% I don't know why, but this works and only includes sections and NOT subsections in the pdf-bookmarks.
\hypersetup{bookmarksdepth=subsection} 

%\usepackage{lmodern}
\usepackage{colortbl}
\usepackage[absolute,overlay]{textpos}
\usepackage{listings}
\usepackage{forloop}
%\usepackage{algorithmic} % PseudoCode package 

\usepackage{tikz}
\usetikzlibrary{matrix}
\usetikzlibrary{arrows.meta}
\usetikzlibrary{automata}
\usetikzlibrary{tikzmark}

% Needed for gbi-macros
\usepackage{xspace}

%%%%%%%%%%%% INHALT %%%%%%%%%%%%%%%%

%% Wochennummer
\newcounter{weeknum}



%% Titelinformationen
\title[GBI-Tutorium, Woche \theweeknum]{Grundbegriffe der Informatik \\ Tutorium \mytutnumber}

\subtitle{Termin \theweeknum \ | \mydate{\theweeknum} \\ \myname}
\author[\myname]{\myname (\mailto{\mymail})}
\institute{KIT -- Karlsruher Institut für Technologie}
\date{\mydate{\theweeknum}\ }

% Modified, DJ (better safe than sorry)
\AuthorTitleSep{ – }

%% Titel einfügen
\newcommand{\titleframe}{\frame{\titlepage}}

%% Alles starten mit \starttut{X}
\newcommand{\starttut}[1]{\setcounter{weeknum}{#1}\titleframe\frame{\frametitle{Inhalt}\tableofcontents} \AtBeginSection[]{%
\begin{frame}
	\tableofcontents[currentsection]
\end{frame}\addtocounter{framenumber}{-1}}}


\newcommand{\framePrevEpisode}{
	\begin{headframe}
		\mylasttimestext
	\end{headframe}
}

%% Legacy: Lastframe. Not for further usage!
\newcommand{\lastframe}[4]{\xkcdframe{#1}{#2}{#3}{#4}{}}
\newcommand{\xkcdframe}[5]{
	\frame[plain]{
		\vspace{-#2pt}
		\begin{figure}[H]
			\centering
			\includegraphics[scale=#1]{#3}
			\vspace{-7pt}
			\caption{ \texttt{\url{#4}} }
			\vspace{5pt}
			#5
		\end{figure} 
	}
}

%% Wörter
\newcommand{\code}[1]{$\mathbf{#1}$}

%% Sterne

\newcounter{starsc}
\newcommand{\stars}[1]{
	\hfill
	\begin{minipage}{100px}
		\forloop{starsc}{0}{\value{starsc} < #1}%
		{%
			\includegraphics[scale=0.05]{star-full.pdf} \hspace*{1px}
		}%
		\forloop{starsc}{\value{starsc}}{\value{starsc} < 5}%
		{%
			\includegraphics[scale=0.05]{star-empty.pdf} \hspace*{1px}
		}
		\vspace*{2px}
	\end{minipage}
}

% No stars for me...
\renewcommand{\stars}[1]{}

\newcommand{\slideThanks}{
	\begin{frame}
		\frametitle{Credits}
		\begin{block}{}
			Vorgänger dieses Foliensatzes wurden erstellt von:\\[1em]
			Thassilo Helmold \\
			Philipp Basler \\
			Nils Braun \\
			Dominik Doerner \\
			Ou Yue \\
		\end{block}
	\end{frame}
}

%% Verbatim
\usepackage{moreverb}